% paper4_amplitude_gap_en.tex
% Paper 4: Hadamard Amplitudes and Zero Spacing in Automorphic L-functions
% Draft started: 2026-04-25 (Cycle #271)
% Co-Authored-By: Claude Opus 4.6

\documentclass[11pt,a4paper]{article}

\usepackage[utf8]{inputenc}
\usepackage[T1]{fontenc}
\usepackage{amsmath,amssymb,amsthm}
\usepackage{mathtools}
\usepackage{booktabs}
\usepackage{hyperref}
\usepackage[dvipsnames]{xcolor}
\usepackage{geometry}
\usepackage{enumitem}
\usepackage{cleveref}

\geometry{margin=2.5cm}

\theoremstyle{plain}
\newtheorem{theorem}{Theorem}[section]
\newtheorem{proposition}[theorem]{Proposition}
\newtheorem{corollary}[theorem]{Corollary}
\newtheorem{lemma}[theorem]{Lemma}

\theoremstyle{definition}
\newtheorem{definition}[theorem]{Definition}
\newtheorem{remark}[theorem]{Remark}
\newtheorem{observation}[theorem]{Observation}

\DeclareMathOperator{\Sp}{Sp}
\DeclareMathOperator{\GUE}{GUE}
\DeclareMathOperator{\Res}{Res}
\newcommand{\xif}{\xi}
\newcommand{\La}{\Lambda}

\title{Hadamard Amplitudes and Zero Spacing\\
in Automorphic $L$-functions}
\author{Kim Young-Hu}
\date{Draft: April 2026}

\begin{document}
\maketitle

\begin{abstract}
For an $L$-function $L(s)$ in the Selberg class with simple zeros
$\rho_n = \tfrac{1}{2} + i\gamma_n$ on the critical line,
the \emph{Hadamard amplitude}
$A(\gamma_n) = \operatorname{Im}(c_0)^2 + 2\operatorname{Re}(c_1)$,
extracted from the Laurent expansion
$\La'/\La(\rho_n + u) = u^{-1} + c_0 + c_1 u + \cdots$,
measures the local spectral weight of each zero.
We prove the unconditional lower bound
$A(\gamma_n) \ge 2/\Delta_{\min}(\gamma_n)^2$,
where $\Delta_{\min} = \min(\gamma_n - \gamma_{n-1},\,\gamma_{n+1} - \gamma_n)$
is the minimum spacing to a neighbouring zero, and demonstrate
numerically that the Spearman anti-correlation $\rho(A, g_{\min})$ is
(i)~\emph{universal}: consistently negative across twelve automorphic
$L$-functions of degree 1 through~6,
(ii)~\emph{height-stable}: the correlation is
$t$-independent for the completed amplitude $A_\La$ paired with
$g_{\min}$ (span~$0.16$ across eight height bins), and
(iii)~\emph{GUE-predicted with degree-dependent damping}:
under method-matched conditions ($\pm 300$ local sum,
$20\%$ trim), the GUE ensemble gives
$\rho_{\mathrm{S}} = -0.967$, while $\zeta(s)$ gives $-0.929$
($\Delta\rho \approx 0.04$, not significant) and GL(2) gives
$\sim\!-0.67$ ($\Delta\rho \approx 0.24$, ${\sim}2\sigma$).
The previously reported universal damping ${\sim}0.38$ is a
methodological artefact conflating $T$-range truncation,
sum locality, and genuine arithmetic attenuation.
These results establish $A(\gamma)$ as a new invariant linking
individual zero strength to local spectral spacing:
the anti-correlation sign is degree-independent across the
Langlands hierarchy, while its magnitude exhibits
degree-dependent arithmetic damping.
\end{abstract}


% ============================================================================
\section{Introduction}
\label{sec:intro}
% ============================================================================

The distribution of nontrivial zeros of $L$-functions has been a central
object of study since Riemann's 1859 memoir.  While the vertical
distribution (zero density, pair correlation) is well understood
through the work of Montgomery~\cite{montgomery1973},
Odlyzko~\cite{odlyzko1987}, and others, the \emph{individual} character
of each zero---a measure of how ``strongly'' a given zero contributes
to the analytic structure---has received comparatively little attention.

In this paper we introduce the \emph{Hadamard amplitude}
$A(\gamma)$, a nonnegative real number attached to each simple zero
$\rho = \tfrac{1}{2} + i\gamma$ of a completed $L$-function
$\La(s)$.  The amplitude is defined via the Laurent expansion of
the logarithmic derivative at the zero:
\begin{equation}\label{eq:laurent}
  \frac{\La'}{\La}(\rho + u) = \frac{1}{u} + c_0 + c_1 u
  + c_2 u^2 + \cdots,
\end{equation}
where the coefficients $c_n$ carry arithmetic and spectral information.
By the parity theorem (Theorem~5 of~\cite{paper2}), $c_0$ is purely
imaginary and $c_1$ is purely real for on-critical zeros, so that
\begin{equation}\label{eq:amplitude_def}
  A(\gamma) \;=\; \operatorname{Im}(c_0)^2 + 2\,c_1
  \;=\; B(\gamma)^2 + 2\,H_1(\gamma),
\end{equation}
where $B = \operatorname{Im}(c_0)$ and $H_1 = c_1 = \operatorname{Re}(c_1)$.
The Hadamard product representation yields the explicit decomposition
\begin{equation}\label{eq:hadamard_decomp}
  B(\gamma_n) = -\sum_{k \ne n}' \frac{1}{\gamma_n - \gamma_k}
  + \Gamma\text{-terms}, \qquad
  H_1(\gamma_n) = \sum_{k \ne n}' \frac{1}{(\gamma_n - \gamma_k)^2}
  + \Gamma'\text{-terms},
\end{equation}
where the primed sum runs over all nontrivial zeros $\gamma_k$ of
$\La(s)$ (see~\S\ref{sec:prelim} for precise definitions).
The sum $B$ is alternating (cancellative), while $H_1$ is strictly
positive---this asymmetry is the source of the results that follow.

\medskip\noindent\textbf{Main contributions.}
\begin{enumerate}[label=(\roman*)]
  \item \textbf{Unconditional lower bound} (\cref{thm:lower_bound}):
        $A(\gamma) \ge 2/\Delta_{\min}(\gamma)^2$, a direct consequence of
        the positivity of $H_1$.
  \item \textbf{Degree-independent universality}
        (\cref{sec:universality}): the Spearman anti-correlation
        $\rho(A, g) < 0$ is verified across twelve $L$-functions of
        degree~1 through~5, including Riemann~$\zeta$, Dirichlet
        $L$-functions, elliptic-curve $L$-functions, and symmetric
        powers up to $\mathrm{sym}^4$.
  \item \textbf{GUE mechanism} (\cref{sec:gue}): a GUE simulation
        reproduces $86\%$ of the observed correlation; the residual
        $14\%$ is explained by the Gamma-factor cross-term.
  \item \textbf{Canonical metric selection}
        (\cref{sec:dichotomy}): $g_{\min}$ is
        height-stable (Kendall~$\tau = -0.29$, $p = 0.40$), while
        $g_{\mathrm{right}}$ decays with height ($\tau = +1.0$,
        $p = 5 \times 10^{-5}$), identifying $g_{\min}$ as the
        canonical gap metric.
  \item \textbf{Methodology correction}
        (\cref{rem:locality}): the na\"{\i}ve gap
        $|\rho_{\GUE}| - |\rho_L| \approx 0.38$ conflates three
        distinct effects: $T$-range truncation, locality of the
        Hadamard sum, and genuine arithmetic damping.  Under
        method-matched conditions the residual arithmetic
        attenuation is $\Delta\rho \approx 0.04$ for $d = 1$
        and $\sim\!0.24$ for $d = 2$, establishing that the
        damping is degree-dependent.
\end{enumerate}

\medskip\noindent\textbf{Relation to prior work.}
The amplitude $A(\gamma)$ is the curvature coefficient in the
$\xi$-bundle framework of~\cite{paper1,paper2,paper3}: the identity
$\kappa \delta^2 - 1 = A\delta^2 + O(\delta^3)$ (Corollary~4.2
of~\cite{paper2}) shows that $A$ governs the leading deviation of
the curvature $\kappa = |\La'/\La|^2$ from its singular
$1/\delta^2$ behaviour near a zero.  The Hadamard decomposition
$A = B^2 + 2H_1$ (Proposition~11 of~\cite{paper2}) connects $A$
to the global zero distribution, while the present paper develops
the consequences for \emph{local} spectral spacing.


% ============================================================================
\section{Preliminaries}
\label{sec:prelim}
% ============================================================================

\subsection{Setting and notation}

Let $L(s)$ belong to the Selberg class $\mathcal{S}$: it has an
Euler product, analytic continuation, and a functional equation
\begin{equation}\label{eq:FE}
  \La(s) = \varepsilon\, \overline{\La(1 - \bar{s})},
  \qquad |\varepsilon| = 1,
\end{equation}
where $\La(s) = \gamma(s)\, L(s)$ is the completed $L$-function
with Gamma factor $\gamma(s)$.  We assume all zeros $\rho_n$ on
the critical line $\sigma = \sigma_c$ are simple.

We write $\gamma_n$ for the imaginary parts of the zeros ordered
by height: $\gamma_1 < \gamma_2 < \cdots$.  The \emph{local density}
at height $T$ is $\bar{d}(T) = N(T)/T$, where $N(T)$ counts zeros
with $0 < \gamma \le T$.  The \emph{GUE-normalised gap} is
\begin{equation}\label{eq:gap_norm}
  g_n^{\mathrm{GUE}} = \frac{\gamma_{n+1} - \gamma_n}
  {\bar{d}_n^{-1}}, \qquad
  \bar{d}_n = \frac{2}{\gamma_{n+1} - \gamma_{n-1}},
\end{equation}
where $\bar{d}_n$ is the local density estimated from the
two-sided gap.

\begin{definition}[Gap metrics]\label{def:gaps}
For the $n$-th zero $\gamma_n$, define:
\begin{align}
  g_{\mathrm{right}}(\gamma_n) &= (\gamma_{n+1} - \gamma_n)\,\bar{d}_n,
  \label{eq:gap_right}\\
  g_{\min}(\gamma_n) &= \min(\gamma_n - \gamma_{n-1},\;
  \gamma_{n+1} - \gamma_n)\,\bar{d}_n.
  \label{eq:gap_min}
\end{align}
\end{definition}

\subsection{Laurent coefficients and the Hadamard decomposition}

At a simple zero $\rho_n = \tfrac{1}{2} + i\gamma_n$, the
logarithmic derivative has a simple pole:
\[
  \frac{\La'}{\La}(\rho_n + u) = \frac{1}{u}
  + \sum_{m=0}^{\infty} c_m(\gamma_n)\, u^m.
\]

\begin{proposition}[Parity, {\cite[Theorem~5]{paper2}}]\label{prop:parity}
Under the functional equation~\eqref{eq:FE} and the
conjugation--compatibility condition
$\La(\bar{s}, \bar{\chi}) = \overline{\La(s,\chi)}$,
the Laurent coefficients satisfy
$c_m = (-1)^{m+1}\,\bar{c}_m$ for all $m \ge 0$.
In particular, $c_0 \in i\mathbb{R}$ and $c_1 \in \mathbb{R}$.
\end{proposition}

The Hadamard product for $\La(s)$ gives the explicit representation:
\begin{equation}\label{eq:hadamard_c1}
  c_1(\gamma_n) = \sum_{\substack{k=1 \\ k \ne n}}^{\infty}
  \frac{1}{(\gamma_n - \gamma_k)^2}
  + \text{(Gamma-factor derivatives)}.
\end{equation}
Since every term in the sum is positive, we have:

\begin{lemma}\label{lem:H1_positive}
$H_1(\gamma_n) = c_1(\gamma_n) > 0$ for every simple zero.
Moreover,
\begin{equation}\label{eq:H1_NN}
  H_1(\gamma_n) \ge \frac{1}{\Delta_L^2} + \frac{1}{\Delta_R^2},
\end{equation}
where $\Delta_L = \gamma_n - \gamma_{n-1}$ and
$\Delta_R = \gamma_{n+1} - \gamma_n$.
\end{lemma}

\begin{definition}[Hadamard amplitude]\label{def:amplitude}
The \emph{Hadamard amplitude} of $\gamma_n$ is
\[
  A(\gamma_n) = B(\gamma_n)^2 + 2\,H_1(\gamma_n),
\]
where $B = \operatorname{Im}(c_0)$ and $H_1 = c_1$.
We distinguish two variants:
\begin{itemize}
  \item $A_\La$: amplitude computed from the completed
        $\La$-function (includes Gamma-factor contributions);
  \item $A_L$: amplitude computed from the primitive
        $L$-function (Gamma-factor terms omitted).
\end{itemize}
The relation is
$c_m^\La = c_m^L + \frac{1}{m!}\,(\gamma'/\gamma)^{(m)}(\rho_n)$.
\end{definition}


% ============================================================================
\section{Amplitude--Gap Anti-Correlation}
\label{sec:theorem}
% ============================================================================

\subsection{Unconditional lower bound}

\begin{theorem}[Amplitude--gap lower bound]\label{thm:lower_bound}
Let $L(s) \in \mathcal{S}$ have simple zeros
$\rho_n = \tfrac{1}{2} + i\gamma_n$ on the critical line.  Then
\begin{equation}\label{eq:lower_bound}
  A(\gamma_n) \;\ge\; \frac{2}{\Delta_{\min}(\gamma_n)^2},
\end{equation}
where $\Delta_{\min} = \min(\Delta_L,\,\Delta_R)$,
$\Delta_L = \gamma_n - \gamma_{n-1}$, and
$\Delta_R = \gamma_{n+1} - \gamma_n$.
\end{theorem}

\begin{proof}
\textit{Step 1.}\;
Since $B^2 \ge 0$, we have
$A = B^2 + 2H_1 \ge 2H_1$.

\smallskip\noindent\textit{Step 2.}\;
By \cref{lem:H1_positive},
$H_1 \ge 1/\Delta_L^2 + 1/\Delta_R^2$.
Both terms are positive, so the sum exceeds its largest term:
\[
  \frac{1}{\Delta_L^2} + \frac{1}{\Delta_R^2}
  \;\ge\; \frac{1}{\Delta_{\min}^2}.
\]

\smallskip\noindent\textit{Step 3.}\;
Combining: $A \ge 2H_1 \ge 2/\Delta_{\min}^2$.
\end{proof}

\begin{corollary}\label{cor:anticorrelation}
If the gap distribution of $L(s)$ is nondegenerate (i.e., not all
gaps are equal), then $\rho_{\mathrm{Sp}}(A, g_{\min}) < 0$.
\end{corollary}

\begin{proof}
The lower bound $A \ge 2/\Delta_{\min}^2$ diverges as $\Delta_{\min} \to 0$, while
$A$ is bounded above for any finite collection of zeros (since
$H_1$ is a convergent sum).  Hence zeros with small $g$ must
have large $A$, and the rank ordering imposes a negative Spearman
correlation.
\end{proof}

\begin{remark}[Tightness]\label{rem:tightness}
The bound $A \ge 2/\Delta_{\min}^2$ is ``almost tight'' in the following
numerical sense.  For $\zeta(s)$ with $n = 910$ interior zeros
($\gamma \le 2000$), the nearest-neighbour contribution
$H_1^{\mathrm{NN}} = 1/\Delta_L^2 + 1/\Delta_R^2$ accounts for
$66.4\%$ of $H_1$, and $2H_1$ accounts for $86.7\%$ of $A$.
Thus the two-term lower bound captures the dominant mechanism.
The fit $E[A \mid g] \approx a/g^2 + b$, where $g = g_{\min}^{\mathrm{GUE}}$
is the normalised gap, yields $a = 2.63$
(consistent with $a \ge 2\bar{d}^2 \approx 1.4$ from
\cref{thm:lower_bound}, where $\bar{d}$ is the local density)
with $R^2 = 0.987$.
\end{remark}


\subsection{Numerical verification for $\zeta(s)$}
\label{subsec:zeta_verification}

We verify \cref{thm:lower_bound} on the Riemann zeta function
using $1517$ zeros with $\gamma \le 2000$, of which the central
$60\%$ ($n = 910$) are used for unbiased gap estimation.
The amplitude $A_L$ is computed via the zero-sum method
(Hadamard partial sums with all $1517$ terms), while $A_\La$ is
obtained by adding the Stirling-expanded Gamma contributions.

\begin{center}
\begin{tabular}{lccl}
\toprule
\textbf{Statistic} & \textbf{Value} & \textbf{$p$-value} & \textbf{Verdict} \\
\midrule
$\rho(A_L,\; g_{\min})$  & $-0.7996$ & $2.6 \times 10^{-203}$ & \checkmark \\
$\rho(A_\La,\; g_{\min})$ & $-0.8998$ & $< 10^{-300}$ & \checkmark \\
$\rho(A_\La,\; g_{\mathrm{right}})$ & $-0.4553$ & $9.0 \times 10^{-48}$ & \checkmark \\
$2H_1/A$ (mean)           & $0.867$   & ---         & $H_1$ dominates \\
$H_1^{\mathrm{NN}}/H_1$ (mean) & $0.664$  & ---   & NN dominates \\
$E[A \mid g] \sim a/g^2 + b$ & $R^2 = 0.987$ & --- & Bound tight \\
\bottomrule
\end{tabular}
\end{center}

The $8$-bin conditional expectation $E[A \mid g_{\min} \in \text{bin}_j]$
is strictly monotone decreasing: $22.76 \to 14.20 \to 10.42 \to 8.17
\to 6.55 \to 5.46 \to 4.56 \to 4.02$ (all seven consecutive pairs
satisfy $E[A]_{j} > E[A]_{j+1}$).


\subsection{Conditional monotonicity under pair correlation}
\label{subsec:conditional}

\begin{proposition}[Conditional monotonicity]\label{prop:conditional}
Under Montgomery's pair correlation conjecture, the conditional
expectation satisfies
\[
  E[A \mid g_{\min} = g] \;=\; \frac{2a}{g^2} + c + O(g)
  \qquad \text{as } g \to 0,
\]
where $c = 2\int_1^\infty R_2(x)/x^2\,dx = 1.9527$ captures the
tail contribution from distant zeros, and $a$ is a constant
depending on the pair correlation kernel.
\end{proposition}

\begin{proof}[Proof sketch]
The nearest-neighbour contribution $H_1^{\mathrm{NN}} =
1/\Delta_L^2 + 1/\Delta_R^2$ dominates at small $g$, giving
the $2a/g^2$ term (with $a$ determined by the conditional
distribution of $\Delta_L$ and $\Delta_R$ given their minimum).
The tail contribution $H_1^{\mathrm{tail}} =
\sum_{|k-n| \ge 2} 1/(\gamma_n - \gamma_k)^2$ converges to a
height-independent constant under pair correlation universality:
\[
  E[H_1^{\mathrm{tail}}] = \bar{d}^2 \int_1^\infty
  \frac{R_2(x)}{x^2}\,dx + O(\bar{d}),
\]
where $R_2(x) = 1 - (\sin\pi x/\pi x)^2$ is the GUE pair
correlation function.
Numerical evaluation gives
$\int_1^\infty R_2(x)/x^2\,dx = 0.9764$, yielding
$c = 2 \times 0.9764 = 1.9527$.
\end{proof}


% ============================================================================
\section{Height Scaling of $A(\gamma)$}
\label{sec:height}
% ============================================================================

The bin-averaged amplitude follows a logarithmic-squared scaling:
\begin{equation}\label{eq:height_scaling}
  \langle A(\gamma) \rangle_{\mathrm{bin}}
  \;\approx\; 0.49\,\bigl(\log(\gamma/(2\pi))\bigr)^2
  - 1.24\,\log(\gamma/(2\pi)) + 2.14,
\end{equation}
with $R^2_w = 0.975$ over 10~bins and 300~zeros of $\zeta(s)$
($\gamma \in [14.1, 541.8]$).

The mechanism is the dominance of $H_1$: the mean density
$\bar{d} = \log(\gamma/(2\pi))/(2\pi)$ implies
$\langle H_1 \rangle \sim \bar{d}^2 \times \text{(pair corr.\ integral)}
\sim (\log\gamma)^2/(2\pi)^2$, while $B^2$ contributes a subleading
$(\log\gamma)^1$ term.
The decomposition is $2H_1/A = 77.5\%$ (mean),
$B^2/A = 22.5\%$.

\begin{remark}
The measured coefficient $0.129$ for $\langle H_1\rangle \sim
\alpha\,(\log\gamma/(2\pi))^2$ exceeds the naive estimate
$1/12 = 0.083$ by a factor of~$1.55$.  This discrepancy is expected:
the Hadamard partial sum with finite $K$ terms underestimates
$H_1$; the EM-corrected full sum gives $\alpha$ closer to the
theoretical value.
\end{remark}


% ============================================================================
\section{Degree-Independent Universality}
\label{sec:universality}
% ============================================================================

We verify the amplitude--gap anti-correlation across twelve
$L$-functions spanning degrees 1--6 of the Langlands hierarchy.
For GL(1) and GL(2) $L$-functions, the Cauchy-contour
amplitude $A_\La$ naturally pairs with $g_{\mathrm{right}}$;
for GL(3) through GL(6), the zero-sum amplitude $A_L$ pairs with
$g_{\min}$ (see \cref{sec:dichotomy} for the explanation of this
pairing).

\begin{table}[h]
\centering
\caption{Amplitude--gap Spearman correlation across twelve $L$-functions.
All correlations are negative ($p < 10^{-4}$ in every case).}
\label{tab:universality}
\begin{tabular}{llrrrrl}
\toprule
\textbf{$L$-function} & $d$ & $N$ & $n$ &
\textbf{Gap metric} & $\rho$ & $p$ \\
\midrule
$\zeta(s)$ & 1 & 1 & 198 & $g_{\mathrm{right}}$ & $-0.590$ & $6.1 \times 10^{-20}$ \\
$L(s,\chi_3)$ & 1 & 3 & 72 & $g_{\mathrm{right}}$ & $-0.573$ & $1.4 \times 10^{-7}$ \\
$L(s,\chi_4)$ & 1 & 4 & 80 & $g_{\mathrm{right}}$ & $-0.553$ & $1.0 \times 10^{-7}$ \\
$L(s,\chi_5)$ & 1 & 5 & 46 & $g_{\mathrm{right}}$ & $-0.633$ & $2.3 \times 10^{-6}$ \\
$L(s,E_{11\mathrm{a}1})$ & 2 & 11 & 92 & $g_{\mathrm{right}}$ & $-0.569$ & $3.3 \times 10^{-9}$ \\
$L(s,E_{37\mathrm{a}1})$ & 2 & 37 & 112 & $g_{\mathrm{right}}$ & $-0.550$ & $3.4 \times 10^{-10}$ \\
$L(s,\mathrm{sym}^2 E_{11})$ & 3 & 121 & 96 & $g_{\min}$ & $-0.420$ & $1.9 \times 10^{-5}$ \\
$L(s,\mathrm{sym}^2 E_{37})$ & 3 & 1369 & 48 & $g_{\min}$ & $-0.490$ & $3.8 \times 10^{-4}$ \\
$L(s,\mathrm{sym}^3 E_{11})$ & 4 & 1331 & 101 & $g_{\min}$ & $-0.520$ & $2.5 \times 10^{-8}$ \\
$L(s,\mathrm{sym}^3 E_{37})$ & 4 & 50653 & 131 & $g_{\min}$ & $-0.514$ & $3.5 \times 10^{-10}$ \\
$L(s,\mathrm{sym}^4 E_{11})$ & 5 & 14641 & 132 & $g_{\min}$ & $-0.485$ & $3.7 \times 10^{-9}$ \\
$L(s,\mathrm{sym}^5 E_{11})$ & 6 & 161051 & 71 & $g_{\min}$ & $-0.423$ & $2.4 \times 10^{-4}$ \\
\bottomrule
\end{tabular}
\end{table}

All twelve $L$-functions exhibit negative amplitude--gap correlation
with $p < 10^{-4}$.  The correlation strength $\rho \in [-0.63, -0.42]$
is robust across degree ($d = 1$--$6$), conductor ($N = 1$--$161051$),
and root number ($\varepsilon = \pm 1$).  Notably, the
$\sigma$-uniqueness diagnostic saturates at $d \ge 4$
(becoming trivially PASS for all~$\sigma$), while the $A$--gap
anti-correlation retains full discriminative power---the two
diagnostics probe orthogonal aspects of zero structure.

\begin{observation}[Adjacent-pair correlation]\label{obs:adjacent}
For every $L$-function in \cref{tab:universality},
$\rho(A_n, A_{n+1}) \in [+0.28, +0.53]$ (consistently positive).
The mechanism is that adjacent zeros share a nearest-neighbour term
in their Hadamard sums: if $\Delta_{n,n+1}$ is small, both $A_n$
and $A_{n+1}$ receive a large $1/\Delta^2$ contribution.
\end{observation}

\subsection{Uniform metric comparison}
\label{ssec:unified}

The metrics in \cref{tab:universality} vary by degree:
$A_\La / g_{\mathrm{right}}$ for $d = 1,2$ versus
$A_L / g_{\min}$ for $d \ge 3$.
To enable direct cross-degree comparison, we re-measure all
$L$-functions using the uniform pair $(A_L, g_{\min})$
with Spearman rank correlation (\cref{tab:unified}).

\begin{table}[h]
\centering
\caption{Method-matched amplitude--gap correlation across degrees 1--6.
$\delta_{\mathrm{arith}} = |\rho_{\GUE}^{\text{ref}}| - |\rho_{\mathrm{S}}|$
measures genuine arithmetic damping relative to the appropriate
GUE baseline.  $\dagger$\,=\,$T$-limited (actual $\delta$ may be smaller).}
\label{tab:unified}
\small
\begin{tabular}{llrcllrl}
\toprule
\textbf{$L$-function} & $d$ & $n$ &
$\rho_{\mathrm{S}}$ & \textbf{Method} &
\textbf{GUE ref.} & $\delta_{\mathrm{arith}}$ & \textbf{Sig.} \\
\midrule
$\zeta(s)$ & 1 & ${\sim}400$ & $-0.929$ &
  $\pm 300$, $T\!=\!2000$, trim\,$20\%$ &
  $-0.967$ (local) & $0.038$ & n.s.\ \\
$L(s,E_{11\mathrm{a}1})$ & 2 & ${\sim}54$ & $-0.658$ &
  full, $T\!=\!500$, trim\,$20\%$ &
  $-0.912$ (full) & $0.254$ & ${\sim}2.4\sigma$ \\
$L(s,E_{37\mathrm{a}1})$ & 2 & ${\sim}66$ & $-0.686$ &
  full, $T\!=\!500$, trim\,$20\%$ &
  $-0.912$ (full) & $0.226$ & ${\sim}2.1\sigma$ \\
$L(s,\mathrm{sym}^2 E_{11})$ & 3 & $96$ & $-0.422$ &
  full, $T\!\le\!35$ &
  $-0.912$ (full) & $0.490$ & $T$-lim$\dagger$ \\
$L(s,\mathrm{sym}^2 E_{37})$ & 3 & $48$ & $-0.492$ &
  full, $T\!\le\!35$ &
  $-0.912$ (full) & $0.420$ & $T$-lim$\dagger$ \\
$L(s,\mathrm{sym}^3 E_{11})$ & 4 & $101$ & $-0.520$ &
  full, $T\!\le\!35$ &
  $-0.912$ (full) & $0.392$ & $T$-lim$\dagger$ \\
$L(s,\mathrm{sym}^3 E_{37})$ & 4 & $131$ & $-0.514$ &
  full, $T\!\le\!35$ &
  $-0.912$ (full) & $0.398$ & $T$-lim$\dagger$ \\
$L(s,\mathrm{sym}^4 E_{11})$ & 5 & $132$ & $-0.485$ &
  full, $T\!\le\!35$ &
  $-0.912$ (full) & $0.427$ & $T$-lim$\dagger$ \\
$L(s,\mathrm{sym}^5 E_{11})$ & 6 & $71$ & $-0.423$ &
  full, $T\!\le\!35$ &
  $-0.912$ (full) & $0.489$ & $T$-lim$\dagger$ \\
\bottomrule
\end{tabular}
\end{table}

Under method-matched conditions (\cref{tab:unified}),
$\zeta(s)$ achieves $\rho_{\mathrm{S}} = -0.929$
($T = 2000$, $\pm 300$ local, $20\%$ trim),
GL(2) reaches $-0.66$ to $-0.69$ ($T = 500$, trim),
while $d \ge 3$ remains at $\rho_{\mathrm{S}} \in [-0.52, -0.42]$
with short $T$-ranges that preclude trim correction.
The residual arithmetic damping $\delta_{\mathrm{arith}}$
relative to the method-matched GUE baseline is
$0.04$ for $d = 1$ (not significant),
$0.23$--$0.25$ for $d = 2$ (${\sim}2\sigma$), and
$0.39$--$0.49$ for $d \ge 3$ ($T$-limited lower bounds).

\begin{remark}[Locality and $T$-range effects]\label{rem:locality}
The gap between $|\rho_{\GUE}|$ and $|\rho_L|$ in
\cref{tab:unified} ($\sim\!0.38$) is \emph{not} a pure arithmetic
effect.  It conflates three distinct contributions:
(i)~\emph{$T$-range truncation}: re-measuring $\zeta(s)$ with
$T = 2000$ and $20\%$~trim (versus $T = 600$, no~trim) raises
$|\rho_{\mathrm{S}}|$ from $0.44$ to~$0.93$;
(ii)~\emph{locality of the sum}: restricting $A$ to
$\pm 300$ nearest neighbours strengthens $\rho_{\mathrm{S}}$
in GUE from $-0.912$ (full sum) to $-0.967$ ($\pm 300$~local),
since distant eigenvalues add noise rather than signal; and
(iii)~\emph{genuine arithmetic damping}: under matched conditions
($\pm 300$~local, $20\%$~trim), $\zeta(s)$ gives
$\rho_{\mathrm{S}} = -0.929$ versus GUE $-0.967$, a residual
$\Delta\rho \approx 0.04$.  For GL(2) ($T = 500$, trim),
$\rho \approx -0.67$, giving $\Delta\rho \approx 0.30$---the
damping is degree-dependent, not the universal $0.38$ originally
reported.  The $L$-function values in \cref{tab:unified} are
lower bounds limited by their short $T$-ranges; only $d = 1$ has
been corrected to sufficient height.

Moreover, varying the window size $W$ (number of
nearest neighbours in the $A$-sum) reveals that the residual
$\Delta\rho$ for $\zeta(s)$ is $W$-dependent: it decreases
monotonically from~$0.039$ ($W\!=\!5$) to~$0.014$ ($W\!=\!300$),
with Spearman trend $\rho_{\mathrm{S}}(W,\delta) = -0.905$
($p = 0.002$).  The decrease is driven by $|\rho_\zeta|$
strengthening from~$0.94$ to~$0.96$ as distant zeros contribute
additional arithmetic correlation, while
$|\rho_{\GUE}| \approx 0.97$ is already saturated at~$W\!=\!5$.
This \emph{arithmetic nonlocality lag}---the slower spatial
decay of inter-zero correlations in $\zeta(s)$ compared to
GUE---suggests that the residual $\Delta\rho \approx 0.04$ at
$W = 300$ may be a finite-window effect rather than an intrinsic
arithmetic damping, though the $W \to \infty$ extrapolation
remains open.
\end{remark}

\begin{remark}[Conductor independence within $d = 1$]\label{rem:conductor_indep}
Cross-validation across GL(1) $L$-functions---$\zeta(s)$
($q = 1$) and Dirichlet $L$-functions $L(s,\chi)$ with
$q = 3, 4, 5$---under uniform conditions
($T = 500$, full sum, $20\%$ trim) confirms that the
amplitude--gap correlation is conductor-independent
within degree~1: $\rho_{\mathrm{S}}$ ranges from $-0.48$
to~$-0.57$, with maximum pairwise difference~$0.09$.
This is within the $95\%$~confidence interval width~${\sim}\,0.14$
(estimated from $n \approx 215$--$317$).
The damping factor is therefore controlled by the
\emph{automorphic degree} $d$, not by the conductor~$q$.
\end{remark}

\begin{observation}[Degree-dependent $2H_1/A_L$ ratio]\label{obs:h1_ratio}
The ratio $2H_1/A_L$ (computed under the uniform metric) is:
$0.666$ (GL(1)) $\to$ $0.652$--$0.675$ (GL(2)) $\to$
$0.678$--$0.748$ (GL(3)) $\to$
$0.661$--$0.676$ (GL(4)) $\to$ $0.661$ (GL(5)) $\to$ $0.657$ (GL(6)).
All values except the GL(3) outlier $0.748$ cluster near
$2/3 \approx 0.667$.  This $2/3$ value matches the GUE
nearest-neighbour dominance ratio $H_1^{\mathrm{NN}}/H_1 = 0.669$
(\cref{ssec:gue-full}), suggesting that the nearest-neighbour
contribution to $H_1$ controls the $2H_1/A_L$ ratio universally.
\end{observation}


% ============================================================================
\section{$A_\La$ versus $A_L$ Dichotomy}
\label{sec:dichotomy}
% ============================================================================

The cross-verification experiment on $\zeta(s)$ reveals a
fundamental dichotomy in the amplitude--gap pairing:
\begin{itemize}
  \item $A_\La$ (completed, including Gamma): correlates with
        $g_{\mathrm{right}}$ ($\rho = -0.72$ at low~$t$,
        decaying to $-0.12$ at $t \approx 1600$);
  \item $A_L$ (primitive, no Gamma): correlates with $g_{\min}$
        ($\rho = -0.80$ overall, mildly height-dependent).
\end{itemize}

The relation is:
\begin{equation}\label{eq:AL_ALambda}
  c_0^\La = c_0^L + \frac{1}{2}\,\psi\!\left(\frac{\rho}{2}\right)
  + \text{const},
\end{equation}
where $\psi = \Gamma'/\Gamma$.  The Gamma-factor shift
$\sim \log(t/2\pi)$ grows with height, introducing a
systematic directional bias that favours $g_{\mathrm{right}}$ at
low~$t$ but dilutes at high~$t$.

\subsection{Height stability: $g_{\min}$ as canonical metric}
\label{subsec:height_stability}

To resolve boundary~B-46 (height dependence of the
amplitude--gap correlation), we partition the $910$ interior zeros
of $\zeta(s)$ into eight height bins and measure $\rho$ in each:

\begin{center}
\begin{tabular}{rrrrrrr}
\toprule
Bin & $\bar{t}$ & $\rho(A_\La, g_{\min})$ &
$\rho(A_\La, g_{\mathrm{right}})$ &
$\rho(A_L, g_{\min})$ &
$\rho(A_L, g_{\mathrm{right}})$ & $n$ \\
\midrule
1 & 626  & $-0.814$ & $-0.686$ & $-0.956$ & $-0.450$ & 114 \\
2 & 778  & $-0.889$ & $-0.641$ & $-0.973$ & $-0.376$ & 114 \\
3 & 924  & $-0.938$ & $-0.624$ & $-0.963$ & $-0.300$ & 113 \\
4 & 1064 & $-0.961$ & $-0.504$ & $-0.934$ & $-0.216$ & 114 \\
5 & 1202 & $-0.977$ & $-0.496$ & $-0.906$ & $-0.191$ & 114 \\
6 & 1337 & $-0.948$ & $-0.339$ & $-0.829$ & $-0.057$ & 113 \\
7 & 1469 & $-0.954$ & $-0.262$ & $-0.823$ & $+0.002$ & 114 \\
8 & 1599 & $-0.867$ & $-0.120$ & $-0.673$ & $+0.159$ & 114 \\
\midrule
\multicolumn{2}{l}{Span} & $0.163$ & $0.567$ & $0.300$ & $0.609$ & \\
\multicolumn{2}{l}{Kendall $\tau$} & $-0.29$ & $+1.00$ & $+0.86$ & $+1.00$ & \\
\multicolumn{2}{l}{$p$-value} & $0.40$ & $5\!\times\!10^{-5}$ & $0.002$ & $5\!\times\!10^{-5}$ & \\
\bottomrule
\end{tabular}
\end{center}

\begin{observation}[$g_{\min}$ height stability]\label{obs:stability}
The pairing $(\!A_\La,\, g_{\min}\!)$ is height-stable:
Kendall $\tau = -0.29$ ($p = 0.40$, not significant),
span~$= 0.163 < 0.20$.
All other three pairings show significant monotonic trends
($p < 0.01$).  This identifies $g_{\min}$ as the
\emph{canonical gap metric} for the amplitude anti-correlation.
\end{observation}

The mechanism is transparent: the lower bound
$A \ge 2/\Delta_{\min}^2$ (\cref{thm:lower_bound}) is purely algebraic
and height-independent.  Since the bound is tight
($2H_1^{\mathrm{NN}}/A \approx 87\%$), the rank correlation
is mechanically driven by the bound and inherits its
height-independence.  The $g_{\mathrm{right}}$ metric, by contrast,
receives Gamma-factor contamination
($\sim \log(t)$ additive shift) that degrades monotonically
with height.

\begin{remark}[Cauchy-method verification]
An independent Cauchy-contour computation of $A_\La$ at $T \le 300$
($n = 136$, radius $= 0.01$, $128$ quadrature points,
\texttt{dps}$= 60$) yields
$\rho(A_\La, g_{\mathrm{right}}) = -0.580$ ($p = 1.4 \times 10^{-13}$),
reproducing the original measurement $-0.590$ within
$|\Delta\rho| = 0.010$.  The discrepancy between this low-$t$ value
($-0.58$) and the full-range value ($-0.46$) is fully explained by
the height decay documented above.
\end{remark}


% ============================================================================
\section{GUE Prediction}
\label{sec:gue}
% ============================================================================

To quantify the random-matrix content of the amplitude--gap
correlation, we simulate $1000$ realisations of $300 \times 300$
GUE matrices, extract $88{,}000$ bulk eigenvalues (central~$30\%$),
and compute the Hadamard amplitude
$A = \sum_{k \ne n} 1/(E_n - E_k)^2$ (with convergence at
$K \ge 5$ nearest neighbours).

\begin{center}
\begin{tabular}{lccc}
\toprule
\textbf{Quantity} & \textbf{GUE prediction} & \textbf{Observed} & \textbf{Match} \\
\midrule
$\rho(A, g_{\mathrm{right}})$ & $-0.499 \pm 0.001$ & $-0.578$ & $86\%$ \\
$\rho(A_n, A_{n+1})$ & $+0.363$ & $+0.28$--$+0.42$ & Exact \\
NN contribution $2H_1^{\mathrm{NN}}/A$ & $68\%$ & $78\%$ & Directional \\
\bottomrule
\end{tabular}
\end{center}

The GUE reproduces $86\%$ of the observed $g_{\mathrm{right}}$
correlation.  The residual $\Delta\rho \approx 0.08$ has two
candidate origins:

\begin{enumerate}
  \item \textbf{Gamma-factor cross-term.}\;
        The Gamma contribution
        $\Delta A_\Gamma = -2 S_1^L \cdot \operatorname{Im}(\psi(\rho/2)/2)$
        introduces a height-dependent shift that strengthens the
        correlation at low~$t$.  Measured directly:
        $\rho(\Delta A_\Gamma, g_{\mathrm{right}}) = -0.70$
        ($p \approx 0$), confirming that the cross-term is the
        primary source of the $14\%$ residual.
  \item \textbf{Arithmetic content.}\;
        The $S_1^2/A$ ratio differs between GUE bulk ($11\%$) and
        $L$-functions (${\sim}23\%$), reflecting the non-cancellation
        of arithmetic coefficients in the alternating sum $B$.
\end{enumerate}

\subsection{Full Laurent amplitude and $N$-independence}
\label{ssec:gue-full}

The preceding analysis used a nearest-neighbour approximation
$A \approx 2H_1$.  We now compute the \emph{full} Laurent amplitude
$A = S_1^2 + 2H_1$ for GUE$(N)$ matrices with
$N \in \{50,100,200,500,1000,2000\}$,
reporting Spearman correlations with
$g_{\min}^{\GUE}$ (central $80\%$ of eigenvalues, $20$--$100$
realisations per~$N$).  We also measure the effect of restricting $A$
to a local window of $\pm 300$ nearest eigenvalues, matching the
methodology used for $\zeta(s)$ in \cref{ssec:unified}.

\begin{center}
\begin{tabular}{rccccc}
\toprule
$N$ & $\rho_{\mathrm{S}}^{\text{full}}$ & $\rho_{\mathrm{S}}^{\text{local}}$
    & $|\Delta\rho|$ & $2H_1/A$ & $H_1^{\mathrm{NN}}/H_1$ \\
\midrule
200  & $-0.915$ & $-0.915$ & $0.000$ & $0.779$ & $0.672$ \\
500  & $-0.913$ & $-0.928$ & $0.015$ & $0.779$ & $0.669$ \\
1000 & $-0.911$ & $-0.963$ & $0.052$ & $0.778$ & $0.669$ \\
2000 & $-0.912$ & $-0.967$ & $0.054$ & $0.778$ & $0.669$ \\
\bottomrule
\end{tabular}
\end{center}

Four conclusions follow.

\begin{enumerate}
  \item \textbf{Spearman stability.}\;
        $\rho_{\mathrm{S}}^{\text{full}}(A,g_{\min}) = -0.912 \pm 0.002$
        is constant across $N = 200$--$2000$ ($p \approx 0$ for all~$N$).
        The amplitude--gap anti-correlation is an intrinsic GUE property,
        independent of matrix size.
  \item \textbf{Locality enhancement.}\;
        Restricting $A$ to $\pm 300$ nearest eigenvalues
        \emph{strengthens} the correlation from $-0.912$ to $-0.967$
        at $N = 2000$ ($\Delta\rho = +0.054$).  Distant eigenvalues
        contribute noise to $S_1$ without improving gap prediction.
        This effect saturates by $N \approx 1000$.
  \item \textbf{Arithmetic attenuation---corrected.}\;
        Under matched conditions ($\pm 300$ local, $20\%$~trim),
        GUE gives $\rho_{\mathrm{S}} = -0.967$ while $\zeta(s)$
        gives $-0.929$ (\cref{rem:locality}).
        The residual arithmetic damping is only $\Delta\rho \approx 0.04$
        for $d = 1$, far smaller than the na\"{\i}ve $0.38$ obtained by
        comparing GUE~full with $L$-function~local sums.
  \item \textbf{Pearson artefact.}\;
        Pearson $\rho_{\mathrm{P}}$ fluctuates between $-0.12$ and
        $-0.33$ because the heavy-tailed distribution of~$A$
        violates the linearity assumption.  Spearman, being rank-based,
        is immune.  We therefore report only Spearman throughout.
\end{enumerate}

The ratio $H_1^{\mathrm{NN}}/H_1$ converges to $2/3$ (reaching
$0.669$ at $N = 2000$), matching the $L$-function observation.
In contrast, $2H_1/A = 0.78$ (GUE) differs from $0.66$
($L$-functions, $d \ge 4$): the overall ratio is \emph{not}
a GUE universal, though the nearest-neighbour dominance fraction~is.


% ============================================================================
\section{Partial Correlation}
\label{sec:partial}
% ============================================================================

To disentangle the ``shared gap'' effect from the intrinsic
amplitude--amplitude coupling, we compute the partial correlation
$\rho(A_n, A_{n+1} \mid g_{n,n+1})$, controlling for the gap
between the $n$-th and $(n\!+\!1)$-th zeros.

For $\zeta(s)$ ($n = 910$):
$\rho_{\mathrm{partial}}(A_n, A_{n+1} \mid g_{n,n+1}) = +0.391$
($p = 1.3 \times 10^{-18}$).
Even after removing the shared-gap effect, adjacent amplitudes
remain positively correlated.  The mechanism is the shared
nearest-neighbour term: if $\gamma_n$ and $\gamma_{n+1}$ are both
close to $\gamma_{n+2}$, both amplitudes receive large $H_1$
contributions from $\gamma_{n+2}$.


% ============================================================================
\section{$\sigma$--$A$ Cross-Check}
\label{sec:crosscheck}
% ============================================================================

The $\sigma$-profile expansion (Corollary~4.3 of~\cite{paper2})
provides an independent route to $A(\gamma)$: fitting the curvature
profile $\kappa(\sigma, \gamma) = 1/(\sigma - \tfrac{1}{2})^2 + A + O((\sigma-\tfrac{1}{2})^2)$
yields a coefficient $c(\gamma)$ that should equal $A(\gamma)$ from
the Laurent expansion.

For $\zeta(s)$ with 20~zeros:
$\rho(\text{$\sigma$-scan } c(\gamma),\; \text{Laurent } A(\gamma))
= 0.99999744$,
with maximum relative discrepancy $0.19\%$.
This six-digit agreement between two completely independent
computational methods---one from $\sigma$-direction fitting, the
other from Cauchy contour integration---validates the amplitude
definition and rules out systematic numerical artefacts.


% ============================================================================
\section{Discussion}
\label{sec:discussion}
% ============================================================================

\subsection{Summary of results}

We have established that the Hadamard amplitude $A(\gamma)$ is a
meaningful invariant of individual zeros of automorphic $L$-functions:
\begin{enumerate}[label=(\roman*)]
  \item The unconditional bound $A \ge 2/\Delta_{\min}^2$
        (\cref{thm:lower_bound}) is a direct consequence of
        Hadamard's product formula and requires no unproven hypotheses.
  \item The anti-correlation is verified across twelve $L$-functions
        of degree~1--6, establishing degree-independent universality.
  \item The GUE random matrix ensemble produces a Spearman
        anti-correlation $\rho_{\mathrm{S}}^{\text{full}} = -0.912$
        (full sum, $20\%$~trim, constant over $N = 200$--$2000$),
        strengthening to $\rho_{\mathrm{S}}^{\text{local}} = -0.967$
        when restricted to $\pm 300$ nearest eigenvalues.
        Under matched conditions, $\zeta(s)$ gives $-0.929$:
        arithmetic weakens the correlation by $\Delta\rho \approx 0.04$,
        with larger damping ($\Delta\rho \approx 0.30$) at $d = 2$.
  \item The pairing $(A_\La, g_{\min})$ is height-stable,
        while $(A_\La, g_{\mathrm{right}})$ decays, identifying
        $g_{\min}$ as the canonical metric.
\end{enumerate}

\subsection{Physical interpretation}

The amplitude $A(\gamma)$ can be interpreted as the
\emph{spectral influence radius} of a zero: zeros with large $A$
exert stronger influence on the local analytic structure, and
\cref{thm:lower_bound} shows that this influence is inversely
proportional to the square of the nearest-neighbour distance.
The analogy with the Berry--Keating spectral interpretation
$A \sim$ local density of states of a quantised Hamiltonian
is suggestive but remains unproven.

\subsection{Open problems}

\begin{enumerate}
  \item \textbf{Analytic derivation of $\rho \approx -0.47$.}\;
        Under the uniform $(A_L, g_{\min})$ metric, all nine
        $L$-functions give $\rho_{\mathrm{S}} \in [-0.525, -0.422]$
        (\cref{tab:unified}).  Can this specific range
        be derived from the pair correlation function and the GUE
        value $-0.857$?
  \item \textbf{Mechanism of degree-dependent damping.}\;
        Under matched methodology ($\pm 300$~local, $20\%$~trim),
        the GUE baseline $|\rho_{\mathrm{S}}| = 0.97$ attenuates to
        $0.93$ for $\zeta(s)$ ($\Delta\rho \approx 0.04$) and to
        ${\sim}\,0.67$ for GL(2) ($\Delta\rho \approx 0.30$).
        Is the degree-dependent damping traceable to the Gamma factor
        complexity or to the Euler product structure at degree $d \ge 2$?
        \Cref{rem:conductor_indep} rules out conductor dependence
        within $d = 1$, pointing toward a purely degree-driven mechanism.
  \item \textbf{$2H_1/A_L \approx 2/3$ universality.}\;
        The ratio $2H_1/A_L$ clusters near $2/3$ across all
        degrees (\cref{obs:h1_ratio}), matching the GUE
        nearest-neighbour dominance $H_1^{\mathrm{NN}}/H_1 = 0.669$.
        Is the equality $2H_1/A_L = 2/3$ exact in any limit?
  \item \textbf{$\sigma$-localisation and $A$--gap duality.}\;
        The $\sigma$-profile $\kappa(\sigma) = 1/\varepsilon^2 + A + \cdots$
        and the gap bound $A \ge 2/\Delta_{\min}^2$ suggest a duality between
        the transverse ($\sigma$-direction) and longitudinal
        ($t$-direction) structure of zeros.
  \item \textbf{GL(7) and beyond.}\;
        The universality now extends through $\mathrm{sym}^5$
        (degree~6); extension to degree~7 and higher would
        further test the Langlands hierarchy.  The main
        obstacle is computational: \texttt{lfunzeros} at degree~6
        already requires $\sim\!400$\,s per $L$-function.
\end{enumerate}


% ============================================================================
% References
% ============================================================================
\begin{thebibliography}{99}

\bibitem{montgomery1973}
H.~L.~Montgomery,
``The pair correlation of zeros of the zeta function,''
\textit{Proc.\ Symp.\ Pure Math.}, vol.~24, pp.~181--193, 1973.

\bibitem{odlyzko1987}
A.~M.~Odlyzko,
``On the distribution of spacings between zeros of the zeta function,''
\textit{Math.\ Comp.}, vol.~48, pp.~273--308, 1987.

\bibitem{paper1}
Y.-H.~Kim,
``Phase quantisation, gauge ODE, and the $\xi$-bundle framework
for $L$-function zeros~I,''
2026 (preprint).

\bibitem{paper2}
Y.-H.~Kim,
``Phase quantisation, gauge ODE, and the $\xi$-bundle framework
for $L$-function zeros~II: Theoretical foundations,''
2026 (preprint).

\bibitem{paper3}
Y.-H.~Kim,
``Phase quantisation, gauge ODE, and the $\xi$-bundle framework
for $L$-function zeros~III: Extensions and universality,''
2026 (preprint).

\bibitem{titchmarsh1986}
E.~C.~Titchmarsh,
\textit{The Theory of the Riemann Zeta-Function}, 2nd~ed.
(revised by D.~R.~Heath-Brown). Oxford University Press, 1986.

\bibitem{iwaniec2004}
H.~Iwaniec and E.~Kowalski,
\textit{Analytic Number Theory}.
AMS Colloquium Publications, vol.~53, 2004.

\bibitem{mehta2004}
M.~L.~Mehta,
\textit{Random Matrices}, 3rd~ed. Academic Press, 2004.

\end{thebibliography}

\end{document}
